\appendixsection{机器人运动学方程完整推导}\label{app:kinematics}

本附录给出第三章中未完全展开的左右腿运动学方程。为便于与控制器状态变量对应，记左、右轮转角为 $\theta_{w,l}$、$\theta_{w,r}$，左、右虚拟腿摆角为 $\theta_{l,l}$、$\theta_{l,r}$，左、右虚拟腿长度为 $l_l$、$l_r$，轮半径为 $R_w$，半轮距为 $R_l$。

\appendixsubsection{位置方程}

\begin{gather}
  s = \frac{R_w(\theta_{w,l}+\theta_{w,r})}{2} \\
  h_b = \frac{l_l\cos\theta_{l,l}}{2}+\frac{l_r\cos\theta_{l,r}}{2} \\
  s_{l,l} = R_w\theta_{w,l}+l_{w,l}\sin\theta_{l,l} \\
  s_{l,r} = R_w\theta_{w,r}+l_{w,r}\sin\theta_{l,r} \\
  h_{l,l} = \frac{l_l\cos\theta_{l,l}}{2}+\frac{l_r\cos\theta_{l,r}}{2}-l_{b,l}\cos\theta_{l,l} \\
  h_{l,r} = \frac{l_l\cos\theta_{l,l}}{2}+\frac{l_r\cos\theta_{l,r}}{2}-l_{b,r}\cos\theta_{l,r} \\
  \phi = \frac{R_w(-\theta_{w,l}+\theta_{w,r})}{2R_l}
  -\frac{l_l\sin\theta_{l,l}}{2R_l}
  +\frac{l_r\sin\theta_{l,r}}{2R_l} \\
  s_b = \frac{R_w(\theta_{w,l}+\theta_{w,r})}{2}
  +\frac{l_l\sin\theta_{l,l}}{2}
  +\frac{l_r\sin\theta_{l,r}}{2}
\end{gather}

\appendixsubsection{速度与加速度方程}

\begin{gather}
  \dot{s} = \frac{R_w\dot{\theta}_{w,l}}{2}+\frac{R_w\dot{\theta}_{w,r}}{2} \\
  \ddot{s} = \frac{R_w\ddot{\theta}_{w,l}}{2}+\frac{R_w\ddot{\theta}_{w,r}}{2} \\
  \dot{\phi} =
  -\frac{R_w\dot{\theta}_{w,l}}{2R_l}
  +\frac{R_w\dot{\theta}_{w,r}}{2R_l}
  -\frac{\dot{\theta}_{l,l}l_l\cos\theta_{l,l}}{2R_l}
  +\frac{\dot{\theta}_{l,r}l_r\cos\theta_{l,r}}{2R_l} \\
  \ddot{\phi} =
  -\frac{R_w\ddot{\theta}_{w,l}}{2R_l}
  +\frac{R_w\ddot{\theta}_{w,r}}{2R_l}
  -\frac{\ddot{\theta}_{l,l}l_l\cos\theta_{l,l}}{2R_l}
  +\frac{\ddot{\theta}_{l,r}l_r\cos\theta_{l,r}}{2R_l}
  +\frac{\dot{\theta}_{l,l}^2l_l\sin\theta_{l,l}}{2R_l}
  -\frac{\dot{\theta}_{l,r}^2l_r\sin\theta_{l,r}}{2R_l} \\
  \dot{s}_b =
  \frac{R_w\dot{\theta}_{w,l}}{2}
  +\frac{R_w\dot{\theta}_{w,r}}{2}
  +\frac{\dot{\theta}_{l,l}l_l\cos\theta_{l,l}}{2}
  +\frac{\dot{\theta}_{l,r}l_r\cos\theta_{l,r}}{2} \\
  \ddot{s}_b =
  \frac{R_w\ddot{\theta}_{w,l}}{2}
  +\frac{R_w\ddot{\theta}_{w,r}}{2}
  +\frac{\ddot{\theta}_{l,l}l_l\cos\theta_{l,l}}{2}
  +\frac{\ddot{\theta}_{l,r}l_r\cos\theta_{l,r}}{2}
  -\frac{\dot{\theta}_{l,l}^2l_l\sin\theta_{l,l}}{2}
  -\frac{\dot{\theta}_{l,r}^2l_r\sin\theta_{l,r}}{2} \\
  \dot{h}_b =
  -\frac{\dot{\theta}_{l,l}l_l\sin\theta_{l,l}}{2}
  -\frac{\dot{\theta}_{l,r}l_r\sin\theta_{l,r}}{2} \\
  \ddot{h}_b =
  -\frac{\ddot{\theta}_{l,l}l_l\sin\theta_{l,l}}{2}
  -\frac{\ddot{\theta}_{l,r}l_r\sin\theta_{l,r}}{2}
  -\frac{\dot{\theta}_{l,l}^2l_l\cos\theta_{l,l}}{2}
  -\frac{\dot{\theta}_{l,r}^2l_r\cos\theta_{l,r}}{2}
\end{gather}

\begin{gather}
  \dot{s}_{l,l} = R_w\dot{\theta}_{w,l}+\dot{\theta}_{l,l}l_{w,l}\cos\theta_{l,l} \\
  \dot{s}_{l,r} = R_w\dot{\theta}_{w,r}+\dot{\theta}_{l,r}l_{w,r}\cos\theta_{l,r} \\
  \ddot{s}_{l,l} = R_w\ddot{\theta}_{w,l}
  +\ddot{\theta}_{l,l}l_{w,l}\cos\theta_{l,l}
  -\dot{\theta}_{l,l}^2l_{w,l}\sin\theta_{l,l} \\
  \ddot{s}_{l,r} = R_w\ddot{\theta}_{w,r}
  +\ddot{\theta}_{l,r}l_{w,r}\cos\theta_{l,r}
  -\dot{\theta}_{l,r}^2l_{w,r}\sin\theta_{l,r}
\end{gather}

\begin{gather}
  \dot{h}_{l,l} =
  \dot{\theta}_{l,l}\left(-\frac{l_l\sin\theta_{l,l}}{2}+l_{b,l}\sin\theta_{l,l}\right)
  -\frac{\dot{\theta}_{l,r}l_r\sin\theta_{l,r}}{2} \\
  \dot{h}_{l,r} =
  -\frac{\dot{\theta}_{l,l}l_l\sin\theta_{l,l}}{2}
  +\dot{\theta}_{l,r}\left(-\frac{l_r\sin\theta_{l,r}}{2}+l_{b,r}\sin\theta_{l,r}\right)
\end{gather}

\begin{equation}
\begin{aligned}
  \ddot{h}_{l,l} &=
  \ddot{\theta}_{l,l}\left(-\frac{l_l\sin\theta_{l,l}}{2}+l_{b,l}\sin\theta_{l,l}\right)
  -\frac{\ddot{\theta}_{l,r}l_r\sin\theta_{l,r}}{2} \\
  &\quad
  +\dot{\theta}_{l,l}^2\left(-\frac{l_l\cos\theta_{l,l}}{2}+l_{b,l}\cos\theta_{l,l}\right)
  -\frac{\dot{\theta}_{l,r}^2l_r\cos\theta_{l,r}}{2}
\end{aligned}
\end{equation}

\begin{equation}
\begin{aligned}
  \ddot{h}_{l,r} &=
  -\frac{\ddot{\theta}_{l,l}l_l\sin\theta_{l,l}}{2}
  +\ddot{\theta}_{l,r}\left(-\frac{l_r\sin\theta_{l,r}}{2}+l_{b,r}\sin\theta_{l,r}\right) \\
  &\quad
  -\frac{\dot{\theta}_{l,l}^2l_l\cos\theta_{l,l}}{2}
  +\dot{\theta}_{l,r}^2\left(-\frac{l_r\cos\theta_{l,r}}{2}+l_{b,r}\cos\theta_{l,r}\right)
\end{aligned}
\end{equation}

\appendixsection{机器人动力学方程完整推导}\label{app:dynamics}

本附录给出第三章线性状态空间模型所依据的原始动力学方程、线性化方程组和矩阵形式。式中 $f_l$、$f_r$ 为左右轮地面切向作用力，$F_{ws}$、$F_{wh}$ 为轮端与腿部之间的水平和竖直约束力，$F_{bs}$、$F_{bh}$ 为腿部与机体之间的水平和竖直约束力。

\appendixsubsection{原始动力学方程}

\begin{gather}
  R_w\ddot{\theta}_{w,l}m_w = -F_{ws,l}+f_l \\
  R_w\ddot{\theta}_{w,r}m_w = -F_{ws,r}+f_r \\
  I_w\ddot{\theta}_{w,l} = -R_wf_l+T_{lw,l} \\
  I_w\ddot{\theta}_{w,r} = -R_wf_r+T_{lw,r} \\
  m_l(R_w\ddot{\theta}_{w,l}+\ddot{\theta}_{l,l}l_{w,l}) = -F_{bs,l}+F_{ws,l} \\
  m_l(R_w\ddot{\theta}_{w,r}+\ddot{\theta}_{l,r}l_{w,r}) = -F_{bs,r}+F_{ws,r} \\
  0 = -F_{bh,l}+F_{wh,l}-gm_l \\
  0 = -F_{bh,r}+F_{wh,r}-gm_l
\end{gather}

\begin{equation}
\begin{aligned}
  I_l\ddot{\theta}_{l,l} &=
  -F_{bs,l}l_{b,l}-F_{ws,l}l_{w,l}+T_{bl,l}-T_{lw,l} \\
  &\quad
  +\theta_{l,l}(F_{bh,l}l_{b,l}+F_{wh,l}l_{w,l})
\end{aligned}
\end{equation}

\begin{equation}
\begin{aligned}
  I_l\ddot{\theta}_{l,r} &=
  -F_{bs,r}l_{b,r}-F_{ws,r}l_{w,r}+T_{bl,r}-T_{lw,r} \\
  &\quad
  +\theta_{l,r}(F_{bh,r}l_{b,r}+F_{wh,r}l_{w,r})
\end{aligned}
\end{equation}

\begin{equation}
  m_b\left(\frac{R_w(\ddot{\theta}_{w,l}+\ddot{\theta}_{w,r})}{2}
  +\frac{\ddot{\theta}_{l,l}l_l}{2}
  +\frac{\ddot{\theta}_{l,r}l_r}{2}\right)
  = F_{bs,l}+F_{bs,r}
\end{equation}

\begin{gather}
  0 = F_{bh,l}+F_{bh,r}-gm_b
\end{gather}

\begin{equation}
\begin{aligned}
  I_b\ddot{\theta}_b &=
  -T_{bl,l}-T_{bl,r} \\
  &\quad
  +\theta_b l_c(F_{bh,l}+F_{bh,r})
  -l_c(F_{bs,l}+F_{bs,r})
\end{aligned}
\end{equation}

\begin{gather}
  I_z\ddot{\phi} = R_l(-f_l+f_r) \\
  F_{wh,l} = F_{wh,r}
\end{gather}

\appendixsubsection{线性化后的动力学方程组}

在直立平衡点附近作小角度线性化，并对约束力进行消元，可得以下五个独立方程：

\begin{align}
  &-T_{bl,l}
  +T_{lw,l}\left(1+\frac{l_{b,l}}{R_w}+\frac{l_{w,l}}{R_w}\right)
  +\ddot{\theta}_{l,l}(I_l-l_{b,l}l_{w,l}m_l) \notag\\
  &\quad
  +\ddot{\theta}_{w,l}\left(
  -\frac{I_wl_{b,l}}{R_w}
  -\frac{I_wl_{w,l}}{R_w}
  -R_wl_{b,l}m_l
  -R_wl_{b,l}m_w
  -R_wl_{w,l}m_w\right) \notag\\
  &\quad
  +\theta_{l,l}\left(
  -\frac{gl_{b,l}m_b}{2}
  -\frac{gl_{w,l}m_b}{2}
  -gl_{w,l}m_l\right)=0 \\
  &-T_{bl,r}
  +T_{lw,r}\left(1+\frac{l_{b,r}}{R_w}+\frac{l_{w,r}}{R_w}\right)
  +\ddot{\theta}_{l,r}(I_l-l_{b,r}l_{w,r}m_l) \notag\\
  &\quad
  +\ddot{\theta}_{w,r}\left(
  -\frac{I_wl_{b,r}}{R_w}
  -\frac{I_wl_{w,r}}{R_w}
  -R_wl_{b,r}m_l
  -R_wl_{b,r}m_w
  -R_wl_{w,r}m_w\right) \notag\\
  &\quad
  +\theta_{l,r}\left(
  -\frac{gl_{b,r}m_b}{2}
  -\frac{gl_{w,r}m_b}{2}
  -gl_{w,r}m_l\right)=0 \\
  &\ddot{\theta}_{l,l}\left(\frac{l_lm_b}{2}+l_{w,l}m_l\right)
  +\ddot{\theta}_{l,r}\left(\frac{l_rm_b}{2}+l_{w,r}m_l\right) \notag\\
  &\quad
  +\ddot{\theta}_{w,l}\left(\frac{I_w}{R_w}+\frac{R_wm_b}{2}+R_wm_l+R_wm_w\right)
  +\ddot{\theta}_{w,r}\left(\frac{I_w}{R_w}+\frac{R_wm_b}{2}+R_wm_l+R_wm_w\right) \notag\\
  &\quad
  -\frac{T_{lw,l}}{R_w}-\frac{T_{lw,r}}{R_w}=0 \\
  &I_b\ddot{\theta}_b+T_{bl,l}+T_{bl,r}
  -\ddot{\theta}_{l,l}l_cl_{w,l}m_l
  -\ddot{\theta}_{l,r}l_cl_{w,r}m_l \notag\\
  &\quad
  +\ddot{\theta}_{w,l}\left(-\frac{I_wl_c}{R_w}-R_wl_cm_l-R_wl_cm_w\right)
  +\ddot{\theta}_{w,r}\left(-\frac{I_wl_c}{R_w}-R_wl_cm_l-R_wl_cm_w\right) \notag\\
  &\quad
  -\theta_bgl_cm_b+\frac{T_{lw,l}l_c}{R_w}+\frac{T_{lw,r}l_c}{R_w}=0 \\
  &-\frac{I_z\ddot{\theta}_{l,l}l_l}{2R_l}
  +\frac{I_z\ddot{\theta}_{l,r}l_r}{2R_l}
  +\frac{R_lT_{lw,l}}{R_w}
  -\frac{R_lT_{lw,r}}{R_w} \notag\\
  &\quad
  +\ddot{\theta}_{w,l}\left(-\frac{I_wR_l}{R_w}-\frac{I_zR_w}{2R_l}\right)
  +\ddot{\theta}_{w,r}\left(\frac{I_wR_l}{R_w}+\frac{I_zR_w}{2R_l}\right)=0
\end{align}

\appendixsubsection{矩阵形式}

令
\begin{equation}
  \ddot{\boldsymbol{q}} =
  \begin{bmatrix}
    \ddot{\theta}_{l,l} & \ddot{\theta}_{l,r} & \ddot{\theta}_b &
    \ddot{\theta}_{w,l} & \ddot{\theta}_{w,r}
  \end{bmatrix}^{\mathsf{T}}
  \quad
  \boldsymbol{q} =
  \begin{bmatrix}
    \theta_{l,l} & \theta_{l,r} & \theta_b
  \end{bmatrix}^{\mathsf{T}}
  \quad
  \boldsymbol{u} =
  \begin{bmatrix}
    T_{lw,l} & T_{lw,r} & T_{bl,l} & T_{bl,r}
  \end{bmatrix}^{\mathsf{T}}
\end{equation}
则线性化系统可写为
\begin{equation}
  \boldsymbol{M}\ddot{\boldsymbol{q}}+\boldsymbol{K}\boldsymbol{q}
  =
  \boldsymbol{E}\boldsymbol{u}
\end{equation}

其中 $\boldsymbol{M}=[M_{ij}]_{5\times5}$ 的非零元素为
\begin{align}
  M_{11} &= I_l-l_{b,l}l_{w,l}m_l \\
  M_{14} &=
  -\frac{I_wl_{b,l}}{R_w}
  -\frac{I_wl_{w,l}}{R_w}
  -R_wl_{b,l}m_l
  -R_wl_{b,l}m_w
  -R_wl_{w,l}m_w \\
  M_{22} &= I_l-l_{b,r}l_{w,r}m_l \\
  M_{25} &=
  -\frac{I_wl_{b,r}}{R_w}
  -\frac{I_wl_{w,r}}{R_w}
  -R_wl_{b,r}m_l
  -R_wl_{b,r}m_w
  -R_wl_{w,r}m_w \\
  M_{31} &= l_{w,l}m_l+\frac{m_b(l_{b,l}+l_{w,l})}{2} \\
  M_{32} &= l_{w,r}m_l+\frac{m_b(l_{b,r}+l_{w,r})}{2} \\
  M_{34} &= M_{35}
  = \frac{I_w}{R_w}+\frac{R_wm_b}{2}+R_wm_l+R_wm_w \\
  M_{41} &= -l_cl_{w,l}m_l \qquad
  M_{42} = -l_cl_{w,r}m_l \qquad
  M_{43} = I_b \\
  M_{44} &= M_{45}
  = -\frac{I_wl_c}{R_w}-R_wl_cm_l-R_wl_cm_w \\
  M_{51} &= -\frac{I_z(l_{b,l}+l_{w,l})}{2R_l} \\
  M_{52} &= \frac{I_z(l_{b,r}+l_{w,r})}{2R_l} \\
  M_{54} &= -\frac{I_wR_l}{R_w}-\frac{I_zR_w}{2R_l} \qquad
  M_{55} = \frac{I_wR_l}{R_w}+\frac{I_zR_w}{2R_l}
\end{align}
其余未列出的 $M_{ij}$ 均为零。

$\boldsymbol{K}=[K_{ij}]_{5\times3}$ 的非零元素为
\begin{align}
  K_{11} &= -\frac{gl_{b,l}m_b}{2}-\frac{gl_{w,l}m_b}{2}-gl_{w,l}m_l \\
  K_{22} &= -\frac{gl_{b,r}m_b}{2}-\frac{gl_{w,r}m_b}{2}-gl_{w,r}m_l \\
  K_{43} &= -gl_cm_b
\end{align}
其余未列出的 $K_{ij}$ 均为零。

$\boldsymbol{E}=[E_{ij}]_{5\times4}$ 的非零元素为
\begin{align}
  E_{11} &= -1-\frac{l_{b,l}}{R_w}-\frac{l_{w,l}}{R_w} &
  E_{13} &= 1 \\
  E_{22} &= -1-\frac{l_{b,r}}{R_w}-\frac{l_{w,r}}{R_w} &
  E_{24} &= 1 \\
  E_{31} &= \frac{1}{R_w} &
  E_{32} &= \frac{1}{R_w} \\
  E_{41} &= -\frac{l_c}{R_w} &
  E_{42} &= -\frac{l_c}{R_w} &
  E_{43} &= -1 &
  E_{44} &= -1 \\
  E_{51} &= -\frac{R_l}{R_w} &
  E_{52} &= \frac{R_l}{R_w}
\end{align}
其余未列出的 $E_{ij}$ 均为零。
